\documentclass[12pt,letterpaper]{report}

% ── Geometry & Layout ──
\usepackage[margin=1in]{geometry}
\usepackage{setspace}
\onehalfspacing
\usepackage[indent=0pt, skip=8pt plus 2pt]{parskip}

% ── Typography ──
\usepackage[T1]{fontenc}
\usepackage[protrusion=true, expansion=false]{microtype}

% ── Math ──
\usepackage{amsmath, amssymb, amsfonts}

% ── Tables ──
\usepackage{booktabs}
\usepackage{tabularx}
\usepackage{colortbl}
\usepackage{array}

% ── Colors ──
\usepackage[table]{xcolor}
\definecolor{darkblue}{HTML}{1A1A2E}
\definecolor{midblue}{HTML}{2B3A4F}
\definecolor{accentblue}{HTML}{3D5A80}
\definecolor{linkblue}{HTML}{1565C0}
\definecolor{proved}{HTML}{E8F5E9}
\definecolor{demonstrated}{HTML}{E3F2FD}
\definecolor{motivated}{HTML}{FFF9C4}
\definecolor{conjectured}{HTML}{FFF3E0}
\definecolor{speculative}{HTML}{FFEBEE}
\definecolor{headershade}{HTML}{2B3A4F}
\definecolor{lightgray}{HTML}{F5F5F5}

% ── Headers & Footers ──
\usepackage{fancyhdr}
\pagestyle{fancy}
\fancyhf{}
\fancyhead[R]{\textit{\small\color{gray} The Quantum Stellarator --- Thesis Document}}
\fancyfoot[C]{\thepage}
\renewcommand{\headrulewidth}{0.4pt}

% ── Section Formatting ──
\usepackage{titlesec}
\titleformat{\chapter}[display]
  {\normalfont\LARGE\bfseries\color{darkblue}}
  {\chaptertitlename\ \thechapter}{16pt}{\LARGE}
\titleformat{\section}
  {\normalfont\Large\bfseries\color{midblue}}
  {\thesection}{1em}{}
\titleformat{\subsection}
  {\normalfont\large\bfseries\color{accentblue}}
  {\thesubsection}{1em}{}
\titlespacing*{\chapter}{0pt}{-20pt}{24pt}

% ── Table of Contents ──
\usepackage{tocloft}
\renewcommand{\cftchapfont}{\bfseries\color{darkblue}}
\renewcommand{\cftsecfont}{\color{midblue}}

% ── Lists ──
\usepackage{enumitem}
\setlist[itemize]{leftmargin=*, itemsep=4pt, parsep=2pt}
\setlist[enumerate]{leftmargin=*, itemsep=4pt, parsep=2pt}

% ── Hyperlinks ──
\usepackage[colorlinks=true, linkcolor=linkblue, urlcolor=linkblue, citecolor=linkblue]{hyperref}

% ── Custom Commands ──
\newcommand{\epistemic}[1]{%
  \par\vspace{4pt}%
  {\small\itshape\color{gray} Epistemic status: #1}%
  \par\vspace{4pt}%
}

\newcommand{\ruledline}{%
  \par\vspace{12pt}%
  {\centering\color{lightgray}\rule{0.4\textwidth}{0.5pt}\par}%
  \vspace{12pt}%
}

% ── New column type for colored header cells ──
\newcolumntype{H}{>{\columncolor{headershade}\color{white}\bfseries\small}l}
\newcolumntype{J}{>{\columncolor{headershade}\color{white}\bfseries\small}p}


% ══════════════════════════════════════════════════════════
% DOCUMENT
% ══════════════════════════════════════════════════════════

\begin{document}

% ── Title Page ──
\thispagestyle{empty}
\begin{titlepage}
\centering
\vspace*{3cm}

{\Huge\bfseries\color{darkblue} THE QUANTUM STELLARATOR\par}
\vspace{0.6cm}
{\Large\color{accentblue} A Formally Verified, Multi-Layer Architecture\\
for Topological Quantum Confinement\par}

\vspace{2cm}
{\color{lightgray}\rule{0.5\textwidth}{0.5pt}\par}
\vspace{1.5cm}

{\large Adrian [Surname]\par}
\vspace{0.3cm}
{\color{gray} Independent Researcher\\
Greater Chicago Area\par}

\vspace{3cm}
{\color{gray} Draft --- March 2026\par}
\vspace{0.4cm}
{\small\itshape\color{gray} Epistemic Status: Architecture document.\\
Individual claims rated per Glassbox Ladder.\par}

\end{titlepage}


% ── Table of Contents ──
\pagenumbering{roman}
\tableofcontents
\clearpage
\pagenumbering{arabic}


% ══════════════════════════════════════════════════════════
% CHAPTER 1: ABSTRACT
% ══════════════════════════════════════════════════════════

\chapter{Abstract}

This document presents the architectural thesis for the Quantum Stellarator: a multi-layer topological quantum confinement device whose design is grounded in formal verification rather than empirical iteration alone. The architecture comprises eight physical layers (0--7), each occupying a distinct mathematical domain, unified by a stellarator integration topos with trefoil knot topology. The formal verification substrate---the Geometry of State (GoS) framework---provides 144 zero-dependency theorems in Lean~4 covering Clifford algebra, the Altland--Zirnbauer tenfold classification, and Majorana zero modes, with zero \texttt{sorry} statements in any L2 classification theorem.

The primary experimental contribution, the Proof of Geometric Topological Chirality (PGTC), demonstrates 70\% phonon suppression in quasiperiodic graphene structures via a Harrison/Keating scaling asymmetry mechanism, establishing that electron hopping scales as $d^{-2}$ while phonon spring constants scale as $d^{-4}$---a separation that enables topologically protected thermal management at the device level.

The epistemic methodology---the Glassbox framework---assigns every claim in this document to one of five tiers (Proved, Demonstrated, Motivated, Conjectured, Speculative) and enforces a Rung Rule prohibiting any claim from depending on a less-certain foundation. This document is itself an exercise in radical epistemic transparency: the architecture is ambitious, the verified core is rigorous, and the boundary between the two is explicitly marked throughout.


% ══════════════════════════════════════════════════════════
% CHAPTER 2: INTRODUCTION
% ══════════════════════════════════════════════════════════

\chapter{Introduction and Motivation}

\section{The Verification Gap in Quantum Device Design}

Contemporary quantum device engineering operates with a significant gap between theoretical proposals and formal guarantees. Device architectures are typically validated through numerical simulation and experimental iteration, with formal methods applied---if at all---after the fact. This thesis inverts that relationship: the Quantum Stellarator is designed from verified mathematical foundations upward, with each physical layer mapped to a specific algebraic or topological domain where machine-checked proofs can establish properties before fabrication.

The core claim is not that formal verification replaces experiment---it cannot---but that the design space can be constrained by proof before it is explored by fabrication. A device whose topological invariants are verified in Lean~4 does not need to discover its symmetry class empirically; it inherits that classification from the proof.

\section{Origins and Intellectual Trajectory}

The Quantum Stellarator project emerged from an unusual trajectory. The author's background spans 20+ years in visual systems design across 14+ industries, including FDA advertising compliance work that required measurement protocols with zero ambiguity---a discipline that proved unexpectedly relevant to formal verification. The project originated from a creative writing endeavor (a science fiction world called AIKO / Quantum Data Ocean) in which a soliloquy about an AI character preceded and motivated the first formal Lean~4 proof. The transition from narrative to mathematics was not a break but a deepening: the same structural intuition that organizes visual systems at scale turns out to organize proof architectures.

The author is self-taught in formal verification, topological physics, and condensed matter physics, having begun systematic study in early 2026. This document is produced without institutional affiliation. The phrase ``the compiler is the credential'' captures the epistemological stance: the GoS repository's 144 theorems compile or they do not, regardless of who wrote them.


% ══════════════════════════════════════════════════════════
% CHAPTER 3: EPISTEMIC METHODOLOGY
% ══════════════════════════════════════════════════════════

\chapter{Epistemic Methodology: The Glassbox Framework}

\section{The Five-Tier Lattice}

Every claim in this thesis is assigned an epistemic status from a five-tier lattice. The tiers, their definitions, and their constraints are presented in Table~\ref{tab:epistemic}.

\begin{table}[ht]
\centering
\caption{The Glassbox five-tier epistemic lattice.}
\label{tab:epistemic}
\small
\begin{tabularx}{\textwidth}{l X c l}
\toprule
\textbf{Tier} & \textbf{Definition} & \textbf{Color} & \textbf{Rung Rule} \\
\midrule
\rowcolor{proved}
Proved & Lean~4 / Cubical Agda theorem with zero \texttt{sorry} & Green & Foundation tier \\
\rowcolor{demonstrated}
Demonstrated & Working code or numerical simulation with convergent results & Blue & $\geq$ Proved deps. \\
\rowcolor{motivated}
Motivated & Supported by established theory; formal proof pending & Yellow & $\geq$ Demonstrated deps. \\
\rowcolor{conjectured}
Conjectured & Plausible hypothesis with partial evidence; not yet proven & Orange & $\geq$ Motivated deps. \\
\rowcolor{speculative}
Speculative & Conceptual leap beyond current evidence; labeled clearly & Red & $\geq$ Conjectured deps. \\
\bottomrule
\end{tabularx}
\end{table}

\section{The Rung Rule}

The Rung Rule is the central constraint of the Glassbox framework: no claim may depend on a foundation of lesser epistemic certainty. A Demonstrated result may depend on Proved foundations. A Motivated argument may depend on Proved or Demonstrated results. But a Proved theorem may never rest on a Conjectured premise---to do so would be to launder uncertainty through formalism.

This rule is enforced at the workflow level through the project navigation flowchart, which gates every work product through an epistemic status check before it can be considered complete. The rule also applies recursively: if a dependency's status is downgraded (e.g., from Demonstrated to Conjectured due to new evidence), all claims that depend on it must be re-evaluated.

\section{Radical Transparency and Retraction Propagation}

The Glassbox framework preserves failure data alongside successes. Retracted claims are not deleted but explicitly marked with their retraction history. The Topological Invariants document, for example, retains the retracted ``singularities as type errors'' framing and is flagged for revision---the retraction itself is part of the epistemic record. This practice is not incidental; it is foundational. A framework that hides its errors teaches less than one that shows them.

The \texttt{sorry} accounting in the GoS repository follows the same principle: six Float ring law \texttt{sorry} statements remain in \texttt{CLHoTT.lean} and are explicitly documented. Zero \texttt{sorry} statements appear in any L2 classification theorem. The distinction matters and is tracked.


% ══════════════════════════════════════════════════════════
% CHAPTER 4: ARCHITECTURE
% ══════════════════════════════════════════════════════════

\chapter{Device Architecture: The Layer Stack}

\section{Architectural Overview}

The Quantum Stellarator is organized as an eight-layer stack (layers 0--7) plus an integration topos (the stellarator itself). Each layer pair occupies a distinct physical and mathematical domain. The layers are not merely stacked physically; they are composed categorically, with each layer's topos mapping to specific verification tools and epistemic constraints. The full layer stack is presented in Table~\ref{tab:layers}.

\begin{table}[ht]
\centering
\caption{The Quantum Stellarator layer stack.}
\label{tab:layers}
\small
\begin{tabularx}{\textwidth}{c l l X}
\toprule
\textbf{Layer} & \textbf{Physical System} & \textbf{Domain} & \textbf{Theoretical Anchor} \\
\midrule
\rowcolor{demonstrated}
0--1 & Nb superconductor + Penrose graphene & Algebraic geometry & Cl(2,0) / Cl(3,0) Clifford algebra \\
\rowcolor{demonstrated}
2--3 & NV centers + carbon + phonon glass & Materials physics & Harrison/Keating scaling asymmetry \\
\rowcolor{conjectured}
4--5 & He-3/He-4 superfluids + Ni-62 wire & SVT + ER=EPR & Superfluid vacuum / analog gravity \\
\rowcolor{motivated}
6--7 & Bismuth medium + OAM lasers & ZX calculus control & OAM encoding \\
\rowcolor{conjectured}
Stellarator & Integration topos & Trefoil knot & All layers compose here \\
\bottomrule
\end{tabularx}
\end{table}

\section{Layers 0--1: Superconducting Substrate}

\epistemic{Motivated. Materials properties are well-established; specific quasicrystal geometry is Conjectured.}

The foundation layers comprise a niobium superconductor substrate patterned with Penrose quasicrystal graphene. The Penrose tiling provides aperiodic long-range order without translational symmetry---a property that proves critical for phonon management in the layers above. The mathematical domain is algebraic geometry, with the Clifford algebras $\mathrm{Cl}(2,0)$ and $\mathrm{Cl}(3,0)$ providing the formal structure verified in the GoS framework.

The quasicrystal geometry at this layer is not decorative; it establishes the scaling asymmetry that the PGTC result exploits. The aperiodic tiling creates a phonon environment fundamentally different from periodic crystals, enabling the thermal management properties documented in Chapter~\ref{ch:pgtc}.

\section{Layers 2--3: Phonon Glass and NV Centers}

\epistemic{Demonstrated (PGTC results). The phonon suppression mechanism is numerically validated; NV center integration is Motivated.}

Layers 2--3 house the core thermal management system and the nitrogen-vacancy (NV) center quantum sensing layer. This is where the PGTC preprint's results live: the demonstration that quasiperiodic disorder in graphene structures produces 70\% phonon suppression through a Harrison/Keating scaling asymmetry. The key insight is dimensional: electron hopping integrals scale as $d^{-2}$ (inverse square of interatomic distance) while phonon spring constants scale as $d^{-4}$ (inverse fourth power). In a quasiperiodic lattice where interatomic distances vary, this differential scaling means phonons are disrupted far more severely than electrons---the lattice acts as a ``phonon glass, electron crystal.''

The NV centers provide quantum sensing capability integrated into the thermal management layer, enabling in-situ measurement of local magnetic fields and temperature gradients during device operation.

\section{Layers 4--5: Superfluid Medium}

\epistemic{Conjectured. The SVT and ER=EPR connections are theoretical; the Ni-62 trefoil wire is Speculative.}

The superfluid layers introduce helium-3 and helium-4 as working media, with a nickel-62 wire wound in a trefoil knot topology providing the magnetic confinement geometry. The theoretical anchors are Superfluid Vacuum Theory (SVT)---the hypothesis that the quantum vacuum behaves as a superfluid---and the ER=EPR conjecture connecting Einstein--Rosen bridges to quantum entanglement.

These layers represent the most speculative physics in the architecture. The connection to analog gravity (Volovik, Unruh, Weinfurtner/\v{S}van\v{c}ara) provides theoretical motivation but not proof. The Glassbox framework explicitly labels this: the superfluid layers are Conjectured at best, with the Ni-62 trefoil wire rated Speculative. No downstream claim in this thesis depends on these layers being physically realizable---they are documented as part of the full architectural vision, not as load-bearing foundations for the verified results.

\section{Layers 6--7: Control Layer}

\epistemic{Motivated. ZX calculus is well-established; application to OAM laser control in this geometry is Conjectured.}

The outermost active layers comprise a bismuth medium and orbital angular momentum (OAM) lasers for quantum state control. The mathematical framework is ZX calculus---a graphical language for quantum computation that maps naturally to the category-theoretic structure of the GoS framework. The OAM encoding provides additional degrees of freedom for state manipulation beyond polarization and frequency.

\section{The Stellarator Integration Topos}

\epistemic{Conjectured as a physical device; Motivated as a mathematical structure.}

The stellarator itself is not a ninth layer but a separate topos---a category-theoretic space in which all eight layers compose. Its geometry is a trefoil knot, chosen because the trefoil is the simplest nontrivial knot and provides the topological constraints needed for stable confinement without the rotational symmetry of a tokamak. The stellarator topos is where the verified properties of individual layers must compose coherently; it is the integration test for the entire architecture.


% ══════════════════════════════════════════════════════════
% CHAPTER 5: PGTC
% ══════════════════════════════════════════════════════════

\chapter{Core Contribution: Proof of Geometric Topological Chirality}
\label{ch:pgtc}

\section{The Harrison/Keating Scaling Asymmetry}

The PGTC's central result rests on a scaling asymmetry first identified in the context of amorphous semiconductors but here applied to quasiperiodic graphene. In Harrison's formulation, the electron hopping integral between atoms at distance $d$ scales as
\begin{equation}
t(d) \propto d^{-2},
\end{equation}
while in Keating's valence force field model, the phonon spring constant scales as
\begin{equation}
k(d) \propto d^{-4}.
\end{equation}
The ratio $k/t \propto d^{-2}$ means that as interatomic distance increases, phonon coupling drops quadratically faster than electron coupling.

In a periodic crystal, this asymmetry is irrelevant---all distances are fixed. In a quasiperiodic structure (Penrose tiling, Fibonacci chain), the distribution of interatomic distances creates a landscape where electron transport remains coherent (the $d^{-2}$ variation is tolerable) while phonon transport is severely disrupted (the $d^{-4}$ variation is not). The result is a structure that conducts electrons like a crystal and scatters phonons like a glass---the phonon glass electron crystal (PGEC) concept, here derived from geometric first principles rather than achieved empirically through doping or alloying.

\section{Numerical Results}

The PGTC preprint reports the quantitative results presented in Table~\ref{tab:pgtc}.

\begin{table}[ht]
\centering
\caption{PGTC numerical results from quasiperiodic graphene simulation.}
\label{tab:pgtc}
\small
\begin{tabularx}{\textwidth}{l l X}
\toprule
\textbf{Quantity} & \textbf{Value} & \textbf{Interpretation} \\
\midrule
\rowcolor{proved}
Phonon suppression & $\kappa_\mathrm{QP}/\kappa_\mathrm{ordered} = 0.30$ & 70\% reduction in thermal conductivity \\
\rowcolor{demonstrated}
Localized phonon modes & 12 modes identified & Quasiperiodic disorder creates localization \\
\rowcolor{proved}
Spectral gap ratio & $58.8\times$ & Topological protection margin \\
\rowcolor{demonstrated}
Winding number & $w = -1$ & Nontrivial topological invariant \\
\rowcolor{proved}
MZM edge localization & 99.7\% & Majorana zero modes at boundaries \\
\rowcolor{demonstrated}
Class transition & BDI $\to$ D & Exchange field breaks time-reversal symmetry \\
\bottomrule
\end{tabularx}
\end{table}

The 70\% phonon suppression ($\kappa_\mathrm{QP}/\kappa_\mathrm{ordered} = 0.30$) is the headline result, but the topological quantities are arguably more significant. The winding number $w = -1$ establishes that the system has a nontrivial topological invariant, meaning its edge properties are protected by topology rather than fine-tuning. The 99.7\% edge localization of Majorana zero modes confirms that the topological protection extends to the boundary states relevant for quantum computation.

\section{Class Transition: BDI to D}

The PGTC demonstrates a symmetry class transition in the Altland--Zirnbauer tenfold classification. The unperturbed system belongs to class BDI (time-reversal symmetric, particle-hole symmetric, chiral). Application of an exchange field breaks time-reversal symmetry, transitioning the system to class D (particle-hole symmetric only). This transition is formally verified in the GoS framework's tenfold way theorems---the L2 classification theorems that carry zero \texttt{sorry} statements.

The spectral gap ratio of $58.8\times$ quantifies the robustness of the topological protection: the gap between topologically protected and unprotected states is nearly 60 times the relevant perturbation scale, indicating that the protected states are stable against realistic levels of disorder and thermal noise.


% ══════════════════════════════════════════════════════════
% CHAPTER 6: FORMAL VERIFICATION
% ══════════════════════════════════════════════════════════

\chapter{Formal Verification Framework}

\section{The Geometry of State (GoS) in Lean 4}

The GoS framework is a zero-Mathlib, single-canvas Lean~4 project that provides the formal verification substrate for the Quantum Stellarator. The deliberate exclusion of Mathlib (Lean's community mathematics library) is a design choice, not a limitation: every theorem in GoS is proved from first principles, with zero external dependencies. This ensures that the proof chain is fully auditable and that no implicit assumptions are inherited from library code.

The current repository contains 144 unique zero-dependency theorems, confirmed through a full repository audit (Architecture v3 had overcounted due to duplicate files; this was corrected). The theorems span Clifford algebra $\mathrm{Cl}(2,0)$ and $\mathrm{Cl}(3,0)$, the Altland--Zirnbauer tenfold classification, topological invariants, Majorana zero modes, and Higher Inductive Topos Theory (HiTT). The verification tool stack is summarized in Table~\ref{tab:verification}.

\begin{table}[ht]
\centering
\caption{Formal verification tool stack.}
\label{tab:verification}
\small
\begin{tabularx}{\textwidth}{l l X l}
\toprule
\textbf{Tool} & \textbf{Framework} & \textbf{Domain} & \textbf{Status} \\
\midrule
\rowcolor{proved}
Lean~4 & Geometry of State & Clifford algebra, AZ classification & 144 theorems; 0 sorry (L2) \\
\rowcolor{demonstrated}
Cubical Agda & TopologicalBridge & HoTT, path types, topology & Port in progress \\
\rowcolor{motivated}
Constructor Theory & Constraint docs & Buildability constraints & Intermediate tier \\
\bottomrule
\end{tabularx}
\end{table}

\section{\texttt{sorry} Accounting}

The GoS repository maintains explicit \texttt{sorry} accounting as part of the Glassbox framework's radical transparency commitment:

\begin{itemize}
  \item \textbf{L2 classification theorems:} Zero \texttt{sorry}. Every theorem in the tenfold way classification compiles without \texttt{sorry} statements.
  \item \textbf{\texttt{CLHoTT.lean}:} Six \texttt{sorry} statements remain, all in Float ring law axioms. These represent a known limitation of Lean~4's native floating-point arithmetic---the ring laws (commutativity, associativity, distributivity) do not hold exactly for IEEE~754 floats. These are explicitly quarantined and no classification theorem depends on them.
  \item \textbf{Total unique theorems:} 144, confirmed by repository audit.
\end{itemize}

\section{The Cubical Agda Bridge}

The \texttt{TopologicalBridge} file initiates a port of GoS results from Lean~4 to Cubical Agda, motivated by the observation that topological properties are more naturally expressed in homotopy type theory (HoTT) than in classical dependent type theory. The Lean~4 framework handles geometric/algebraic structures well (Clifford algebras, matrix representations); the Cubical Agda framework handles topological/contextual structures (path types, higher inductive types, univalence). The two tools address complementary domains rather than competing.

\section{The Omnivalence Postulate}

\epistemic{Conjectured.}

The Omnivalence Postulate is a meta-level principle extending Voevodsky's univalence axiom. Where univalence asserts that equivalent types are identical (within HoTT), Omnivalence proposes a bridge principle that allows results to transfer between dependent type theory (Lean~4) and cubical type theory (Cubical Agda) while preserving their epistemic status. This is currently Conjectured---no formal proof exists---and is documented here as a research direction rather than a foundation. No verified result in this thesis depends on the Omnivalence Postulate.


% ══════════════════════════════════════════════════════════
% CHAPTER 7: THEORETICAL FOUNDATIONS
% ══════════════════════════════════════════════════════════

\chapter{Theoretical Foundations}

\section{The Tenfold Way and Cayley--Dickson Ladder}

\epistemic{Motivated.}

The Altland--Zirnbauer tenfold classification of topological insulators and superconductors provides the symmetry framework for the Quantum Stellarator. The GoS framework formalizes all ten symmetry classes and their topological invariants. A deeper structural observation maps the tenfold way to the Cayley--Dickson construction---the algebraic ladder from real numbers through complex numbers, quaternions, octonions, and beyond. Each step in the Cayley--Dickson ladder doubles the algebra's dimension while losing an algebraic property (commutativity, then associativity, then alternativity), and each symmetry class in the tenfold way corresponds to a position on this ladder.

This mapping is Motivated by structural analogy and partially verified through the GoS theorems, but a complete formal proof of the correspondence remains open work.

\section{Hamiltonian Sculpting}

\epistemic{Conjectured.}

Hamiltonian Sculpting is the design methodology implied by the formal verification approach: rather than discovering a system's Hamiltonian through experiment and fitting it to a symmetry class post hoc, one begins with a target symmetry class (from the tenfold way), derives the Hamiltonian constraints from the topos structure, and constructs a physical system that satisfies those constraints by design. The tenfold topoi serve as geometric templates for Hamiltonian construction.

This methodology is Conjectured as a general principle. The PGTC result demonstrates one instance (constructing a BDI system and transitioning it to class D by design), but the generalization to arbitrary symmetry classes and layer compositions has not been formally established.

\section{Cohomological Obstruction Framework}

\epistemic{Motivated.}

The cohomological obstruction framework connects phase transitions to sheaf cohomology's failure-to-glue. In this framing, a topological phase transition occurs when local data (defined on open sets of parameter space) can no longer be consistently glued into a global section---the obstruction to gluing lives in a cohomology group, and the phase transition is the point where that obstruction becomes nontrivial. This provides a unified language for understanding both the BDI$\to$D class transition in the PGTC and the broader landscape of topological phase transitions across the layer stack.

\section{Constitutional XAI}

\epistemic{Conjectured. Identified as a significant unnamed gap in AI alignment research.}

An emergent conceptual contribution from the GoS framework is Constitutional XAI: the idea that AI constitutional principles (behavioral constraints, safety properties, alignment objectives) can be encoded as compiler-enforced type constraints rather than as natural-language rules enforced through RLHF. In this framing, a ``constitutional violation'' becomes a type error caught at compile time rather than a behavioral failure caught at evaluation time. The GoS framework's zero-\texttt{sorry}, zero-dependency architecture provides a proof-of-concept for this approach: the same methodology that verifies topological invariants in physics can verify behavioral invariants in AI systems.


% ══════════════════════════════════════════════════════════
% CHAPTER 8: ENGINEERING ROADMAP
% ══════════════════════════════════════════════════════════

\chapter{Engineering Roadmap and Simulation Strategy}

\section{Verification-First Development Cycle}

The Quantum Stellarator follows a verification-first development cycle in which formal proofs precede simulation and simulation precedes fabrication. The project navigation flowchart encodes this as a decision tree: at every branch point, the developer must assess whether the current work product is formally verified, buildable via Constructor Theory constraints, or still conjectural. The flowchart routes each case to the appropriate next action and gates all outputs through the Glassbox epistemic status check.

\section{Simulation Platforms}

Two simulation platforms are designated for different domains:

\begin{itemize}
  \item \textbf{Azure Quantum:} Materials-level simulation---crystal structures, phonon dispersions, electronic band structure. Appropriate for layers 0--3 where the physics is closest to standard condensed matter.
  \item \textbf{CUDA-Q / Kiskade:} Quantum circuit simulation on NVIDIA GPU backend. Appropriate for layers 6--7 (ZX calculus control) and for the stellarator integration topos where quantum circuit representations are natural.
\end{itemize}

The choice between platforms is domain-driven, not preference-driven. Materials questions go to Azure Quantum; circuit questions go to CUDA-Q. The formal verification layer (GoS/Cubical Agda) sits above both and constrains what is simulated to what has been at least Motivated on the Glassbox ladder.

\section{Targets and Next Steps}

The immediate engineering priorities are:

\begin{enumerate}
  \item \textbf{PGTC arXiv submission:} The \LaTeX{} draft (REVTeX~4-2 format) is complete. Submission is gated on arXiv endorsement, currently being pursued through a postdoctoral connection at the University of Wisconsin.
  \item \textbf{TopologicalBridge completion:} Port remaining GoS topological theorems to Cubical Agda to validate the dual-tool verification strategy.
  \item \textbf{Topological Invariants revision:} Revise the document to remove the retracted ``singularities as type errors'' framing---flagged as highest priority.
  \item \textbf{Industry outreach:} Strategic briefs targeting Medline (formal verification for software-driven supply chains) and Abbott (quality and verification) through family network connections. Broader targets include Microsoft Station Q, NVIDIA Quantum, and research partnerships with Anthropic.
\end{enumerate}


% ══════════════════════════════════════════════════════════
% CHAPTER 9: DISCUSSION
% ══════════════════════════════════════════════════════════

\chapter{Discussion: Scope, Limitations, and Honest Assessment}

\section{What This Thesis Claims}

This thesis claims that a multi-layer topological quantum confinement device can be designed from formally verified mathematical foundations, that the Geometry of State framework provides a working proof-of-concept for this approach, and that the PGTC result demonstrates a specific, numerically validated instance of verification-first physics (phonon suppression from geometric first principles). It claims that the Glassbox framework provides a rigorous epistemic methodology for managing the certainty gradient inherent in any project that spans proved theorems and speculative physics.

\section{What This Thesis Does Not Claim}

This thesis does not claim that the Quantum Stellarator is buildable with current technology. Layers 4--7 remain Conjectured or Speculative, and the stellarator integration topos is a mathematical structure, not an engineering blueprint. The superfluid vacuum theory connections (layers 4--5) are theoretical and may prove physically unrealizable. The Omnivalence Postulate is unproved. The Ni-62 trefoil wire is Speculative.

The thesis also does not claim that formal verification eliminates the need for experiment. It claims that verification can constrain the design space before experiment explores it---a weaker but more defensible position.

\section{The Submission Pattern}

A documented pattern in this project's history is the production of artifacts at high volume and quality followed by a consistent stop before the submission step. The PGTC preprint has been draft-complete for some time; the arXiv endorsement outreach has been identified as the critical path and has not yet been initiated. This pattern is acknowledged here not as self-criticism but as data: the Glassbox framework's commitment to radical transparency extends to the researcher's own behavioral patterns. The flowchart encodes this as an explicit gate (``Text Wisconsin cousin NOW'') precisely because encoding it is easier than doing it. The gap between the two is the current rate-limiting factor for the project's public-facing progress.


% ══════════════════════════════════════════════════════════
% CHAPTER 10: CONCLUSION
% ══════════════════════════════════════════════════════════

\chapter{Conclusion}

The Quantum Stellarator project demonstrates that formal verification and topological physics can be productively unified in a single architectural framework. The Geometry of State provides 144 machine-checked theorems grounding the device's symmetry classification. The PGTC provides a numerically validated physical mechanism (phonon suppression via scaling asymmetry) that emerges from the geometry itself. The Glassbox framework provides the epistemic discipline to distinguish what is proved from what is hoped.

The architecture is ambitious by design---eight layers, multiple mathematical domains, a trefoil stellarator topos---and the epistemic methodology is honest by design. Not every layer is verified. Not every connection is proved. But the boundary between the verified and the conjectural is explicitly marked, the \texttt{sorry} count is public, the retraction history is preserved, and the compiler accepts or rejects without regard to the author's credentials.

The next step is the simplest and the hardest: submit.

\ruledline

\vfill
\begin{center}
{\small\itshape\color{gray} End of Document\\
Draft --- March 2026 --- Epistemic status annotations throughout}
\end{center}

\end{document}
